\section*{الفصل \ref{ch:b3:molecular-evolution} --- التطور الجزيئي والجينوميات التطورية}

\begin{solution}{exo:b3:molecular-evolution:1}
$k = \mu$: أي إن \omterm{thm:b3:molecular-evolution:neutral}{معدل الاستبدال} المحايد لكل موقع في الجيل
يساوي معدل الطفر. ففي كل جيل تنشأ $2N\mu$ طفرة محايدة جديدة
ولكل واحدة احتمال $1/2N$ للتثبّت؛ والجمهرة الأكبر تصنع
طفرات أكثر وتعطي كل واحدة حظًا أصغر بالتناسب، فيلغي الأثران
أحدهما الآخر إلغاءً تامًا.
\end{solution}

\begin{solution}{exo:b3:molecular-evolution:2}
النظائر المباشرة: مورّثات في نوعين منحدرة من مورّثة واحدة في
سلفهما المشترك --- كغلوبين $\beta$ في الإنسان وغلوبين $\beta$ في الفأر.
والنظائر المتضاعفة: مورّثات في جينوم واحد منحدرة من تضاعف ---
كغلوبين $\alpha$ وغلوبين $\beta$ في الإنسان، أو $\beta$ و$\gamma$. \omterm{def:b3:molecular-evolution:duplication}{والمورّثة الكاذبة}:
نسخة مضمحلة لم تعد ترمّز بروتينًا --- كالمورّثة $\psi\beta$
في عنقود غلوبين $\beta$ البشري، بما فيها من رامزات وقف.
\end{solution}

\begin{solution}{exo:b3:molecular-evolution:3}
تقارن $\omega$ \omterm{thm:b3:molecular-evolution:neutral}{معدل الاستبدال} المغيِّر للحمض الأميني \omterm{thm:b3:molecular-evolution:neutral}{بمعدل
الاستبدال} الصامت، وهو خط الأساس المحايد. فالقيمة $0.02$:
انتخاب منقٍّ قوي، إذ تُزال \qty{98}{\%} من تغيّرات الأحماض الأمينية
--- كالأكتين، \omterm{def:b3:chromatin-epigenetics:nucleosome}{والهستون} H3. و$1.0$: بلا قيد --- \omterm{def:b3:molecular-evolution:duplication}{كمورّثة كاذبة}. و$2.5$:
\omterm{def:b3:molecular-evolution:dnds}{انتخاب موجب}، إذ تُحابى تغيّرات الأحماض الأمينية --- كأخدود
ربط الببتيد في MHC، \omterm{def:b3:virology:virus}{وغلاف} HIV، وليزوزيم المجترّات.
\end{solution}

\begin{solution}{exo:b3:molecular-evolution:4}
البروتينات المقارَنة عبر الثدييات تتغير بنحو استبدال واحد لكل
موقع في كل مليار سنة؛ وعلى جينوم من مليار موقع مرمِّز وجيل
من بضع سنين يكون ذلك من رتبة استبدال واحد
مثبَّت في الجيل، أو واحد كل سنتين. وكلفة الانتخاب
عند هالدين تسمح بنحو استبدال منتخَب واحد لكل ثلاثمئة
جيل. فالمعدلان يختلفان بعامل مئة أو أكثر، فوجب أن تكون
كل التغيّرات المثبَّتة تقريبًا محايدة.
\end{solution}

\begin{solution}{exo:b3:molecular-evolution:5}
المحايد: $1/2N = 10^{-4}$. وعند $s = 0.001$، $4Ns = 20$: $(1 - \mathrm{e}^{-0.002})/
(1 - \mathrm{e}^{-20}) = 0.002$. وعند $s = 0.01$: $0.0198$. وعند $s = -0.001$: $(1 -
\mathrm{e}^{0.002})/(1 - \mathrm{e}^{20}) = 0.002\,\mathrm{e}^{-20} =
4\times 10^{-12}$. ولا شيء منها محايد فعليًا: فذلك كان يقتضي $|4Ns|
\lesssim 1$، أي $|s| \lesssim 5\times 10^{-5}$؛ فمعامل انتخاب
قدره عُشر في المئة حاسم أصلًا في جمهرة
من خمسة آلاف.
\end{solution}

\begin{solution}{exo:b3:molecular-evolution:6}
بلا تصحيح: $T = 0.12/(2\times 2\times 10^{-9}) = \qty{30}{Myr}$.
وبجوكس--كانتور: $d = -\tfrac{3}{4}\ln(1 - 0.16) = 0.131$، و$T =
\qty{33}{Myr}$. ويبلغ التصحيح عاملًا قدره اثنان حين $-\tfrac{3}{4}
\ln(1 - 4p/3) = 2p$، أي عند $p \approx 0.6$: وعندئذ يكون الفرق الخام قد
فقد معظم معلومته.
\end{solution}

\begin{solution}{exo:b3:molecular-evolution:7}
$p_{N} = 7/700 = 0.010$، و$p_{S} = 10/200 = 0.050$، و$\omega = 0.2$:
أي انتخاب منقٍّ يزيل أربعة أخماس تغيّرات الأحماض الأمينية. وعند
$35$ و$10$: $p_{N} = 0.05$، و$\omega = 1$: أي لا قيد قابلًا للكشف ---
\omterm{def:b3:molecular-evolution:duplication}{مورّثة كاذبة}، أو مورّثة يوازن فيها \omterm{def:b3:molecular-evolution:dnds}{انتخاب موجب} في بعض
المواقع انتخابًا منقيًا في غيرها.
\end{solution}

\begin{solution}{exo:b3:molecular-evolution:8}
$k = d/2T = 0.012/(1.3\times 10^{7}) = 9.2\times 10^{-10}$ لكل موقع في
السنة؛ و$\mu = 25k = 2.3\times 10^{-8}$ لكل موقع في الجيل؛ و$6\times
10^{9}\times 2.3\times 10^{-8} \approx 140$ طفرة جديدة في الجينوم
الثنائي الصيغة في الجيل. والتسلسل المباشر للوالدين والأبناء يجد
نحو سبعين: فتقدير التباعد منتفخ بتعدد الأشكال الحاضر أصلًا
في الجمهرة السلفية، وهو يضيف بضع مئات
من آلاف الأجيال إلى الانفصال الظاهر.
\end{solution}

\begin{solution}{exo:b3:molecular-evolution:9}
$\mathrm{e}^{-T/2N} = \mathrm{e}^{-1} = 0.37$؛ والنسبة المخالِفة $\tfrac{2}{3}
\times 0.37 = 0.25$. ولأجل \qty{10}{\%}: $\mathrm{e}^{-T/2N} = 0.15$، و$T =
2N\ln(1/0.15) = \num{190000}$ جيل.
\end{solution}

\begin{solution}{exo:b3:molecular-evolution:10}
حين $s \to 0$: $1 - \mathrm{e}^{-2s} \approx 2s$ و$1 - \mathrm{e}^{-4Ns}
\approx 4Ns$، فالنسبة $1/2N$. وعند $4Ns \gg 1$ يكون المقام $1$ والبسط
$\approx 2s$. وعند $s < 0$ مع $4N|s| \gg 1$: البسط $1 -
\mathrm{e}^{2|s|} \approx -2|s|$، والمقام $1 - \mathrm{e}^{4N|s|} \approx
-\mathrm{e}^{4N|s|}$، فيكون $P \approx 2|s|\,\mathrm{e}^{-4N|s|}$. وأما مئة ضعف
دون المحايد عند $N = 10^{4}$: فإن $2|s|\,\mathrm{e}^{-4\times 10^{4}|s|} =
5\times 10^{-7}$؛ و$|s| = 1.6\times 10^{-4}$ يعطي $3.2\times 10^{-4}\times
\mathrm{e}^{-6.4} = 5.3\times 10^{-7}$. فيكون $|s| \approx 1.6\times 10^{-4}$،
و$4N|s| \approx 6.5$: أي إن مساوئ قدرها جزء من مئة من في المئة
تكفي، في عشرة آلاف فرد، لجعل التثبّت أندر مئة مرة
من المصادفة.
\end{solution}

\begin{solution}{exo:b3:molecular-evolution:11}
مع نقطة الانفصال على مسافة $a$ من $O$ على المسارين، $a + m =
0.30$، و$a + h = 0.24$، و$m + h = 0.20$: فيكون $m - h = 0.06$، ومن ثم $m = 0.13$، و$h
= 0.07$، والنسبة $1.9$. والساعة الصارمة في كل جيل كانت ستعطي
نسبة عدّتي الأجيال، أي $25/0.5 = 50$. والمرصود $1.9$ يعني أن
الفأر لم يراكم إلا ضعف تغيّر الإنسان في السنة رغم
خمسين ضعفًا من الأجيال: فلا بد أن يكون معدل الطفر في الجيل
أعلى في الإنسان بنحو $25$ ضعفًا --- أي انقسامات خلوية أكثر في
السلالة الجنسية في الجيل --- حتى لا تختلف السلالتان في السنة
إلا بعامل اثنين.
\end{solution}

\begin{solution}{exo:b3:molecular-evolution:12}
تصل الاستبدالات الصانعة للوقف عند $1000\times 2\times 10^{-9}\times
0.04 = 8\times 10^{-8}$ في السنة: فتُتوقَّع الأولى بعد نحو
\qty{12}{Myr} (وانزياحات الإطار تنصّف ذلك تقريبًا). وفي تلك
النافذة تكون الطفرة النافعة، وهي واحدة من $10^{3}$ إلى $10^{5}$ من كل الطفرات، أقل
احتمالًا بكثير للوصول من رامزة وقف، وحتى حين تصل
تتثبت باحتمال $2s$ فقط. فالنسخة تموت عادةً قبل أن
يجد لها الانتخاب استعمالًا --- وعمر النصف المرصود للنسخ
المضاعَفة في الفقاريات بضعة ملايين من السنين.
\end{solution}

\begin{solution}{pb:b3:molecular-evolution:1}
\textbf{1.} $k = 0.012/(2\times 6.5\times 10^{6}) = 9.2\times 10^{-10}$ لكل
موقع في السنة؛ و$\mu = 25k = 2.3\times 10^{-8}$ لكل موقع في الجيل.
\textbf{2.} $6\times 10^{9}\times 2.3\times 10^{-8} \approx 140$ طفرة
جديدة في الجينوم الثنائي الصيغة؛ ومنها \qty{1.5}{\%}، أي نحو $2$، في
متوالية مرمِّزة.
\textbf{3.} في الجينوم الفردي الصيغة، $3\times 10^{9}\times 2.3\times 10^{-8}
\approx 70$ استبدالًا مثبَّتًا في الجيل، في مقابل $1/300$ عند
هالدين: أي أكثر عشرين ألف مرة مما كان الانتخاب يقدر على
سوقه، فتكون الأغلبية الساحقة محايدة.
\textbf{4.} $2N\mu = 10^{5}\times 2.3\times 10^{-8} = 2.3\times 10^{-3}$
طفرة جديدة لكل موقع في الجيل في الجمهرة؛ وستتثبت منها $1/2N =
10^{-5}$.
\textbf{5.} $4N = \num{200000}$ جيل، أي \qty{5}{Myr} --- أي قرابة
طول الزمن منذ الانفصال.
\textbf{6.} التباعد يعدّ الطفرات التي نشأت على السلالتين
المنفصلتين منذ سلفهما المشترك؛ وكل واحدة تتثبت في سلالتها
هي، والمعدل على المدى الطويل $\mu$ مهما يكن الزمن الذي تستغرقه
كل واحدة، إذ تنشأ الطفرات باطراد ولا يفعل التأخر إلا أن
يؤخر التدفق بلا أن يغيّره. (والتصحيح الوحيد هو لتعدد الأشكال
الحاضر عند الانفصال، وهو يجعل السلف المشترك لأليلين أقدم من
انفصال النوعين.)
\textbf{7.} المحايد $10^{-5}$؛ وعند $s = 0.001$: $0.002$؛ وعند $s = 0.01$: $0.0198$؛
وعند $s = 0.1$: $1 - \mathrm{e}^{-0.2} = 0.18$.
\textbf{8.} $s = 1/4N = 5\times 10^{-6}$: فدون ذلك يغرق الانجرافُ
الانتخابَ وتسلك الطفرة سلوك المحايدة؛ وفوقه يحكم الانتخاب
الحظوظ.
\textbf{9.} عند $s = -0.001$: $P = 0.002\,\mathrm{e}^{-200}$، أي صفر لكل
غرض --- أي $10^{-85}$ من المحايد. وعند $s = -10^{-5}$، $4N|s| = 2$: $P =
2\times 10^{-5}/(\mathrm{e}^{2} - 1) = 3.1\times 10^{-6}$، أي \qty{31}{\%} من
المحايد: فالطفرات الضارة ضررًا خفيفًا تتثبت فعلًا.
\textbf{10.} الاستبدالات النافعة لكل موقع في الجيل: $2N\times
10^{-5}\mu\times 2s = 10^{5}\times 10^{-5}\times 0.02\,\mu = 0.02\mu$؛
والمحايدة: $\approx\mu$. فالنسبة التكيّفية $0.02/1.02 \approx \qty{2}{\%}$.
ومع أن كل طفرة نافعة أرجح تثبّتًا بألفي مرة
من المحايدة، فهي نادرة جدًا حتى لا يكون تكيّفيًا إلا استبدال
واحد من خمسين: فالانتخاب يعمل ولا يكاد يظهر.
\textbf{11.} $p_{N} = 46/920 = 0.050$، و$p_{S} = 84/280 = 0.30$. وبالتصحيح:
$d_{N} = -\tfrac{3}{4}\ln(1 - 0.0667) = 0.052$، و$d_{S} = -\tfrac{3}{4}
\ln(1 - 0.40) = 0.38$. فيكون $\omega = 0.14$.
\textbf{12.} انتخاب منقٍّ يزيل نحو \qty{86}{\%} من
تغيّرات الأحماض الأمينية --- أي مورّثة نموذجية. و$k = d_{S}/2T = 0.38/(1.8\times
10^{8}) = 2.1\times 10^{-9}$ لكل موقع في السنة: أي ضعف قيمة الإنسان
والشمبانزي، كما هو متوقع لمقارنة تشمل سلالة القارض السريعة
الدقّ؛ وهي من رتبة المقدار نفسها.
\textbf{13.} $d = -\ln(1 - 0.55) = 0.80$ لكل موقع؛ و$T = 0.80/(2\times
10^{-9}) = \qty{400}{Myr}$.
\textbf{14.} سلسلتان مختلفتان تصنعان الرباعي $\alpha_{2}\beta_{2}$،
الذي يتوقف ربطه التعاوني للأكسجين \omterm{prop:b3:structural-biology:mechanism}{وأثر بور} وتحكمه الأستيري
بثنائي فوسفوغليسيرات-2,3 على السطح البيني بين وحيدات غير
متشابهة؛ وقد أمكن لعنقود $\beta$ بعد ذلك أن يتنوع إلى سلاسل
جنينية وحملية وبالغة بألفات مختلفة.
\textbf{15.} $2kT = 2\times 8\times 10^{-9}\times 4\times 10^{7} = 0.64$
استبدال لكل موقع؛ ومع التشبع بين عشرين حمضًا أمينيًا يكون
الفرق الخام نحو $0.95(1 - \mathrm{e}^{-0.64/0.95}) \approx 0.47$.
وعند \qty{500}{Myr}، $d = 8$: فيكون كل موقع قد تغيّر مرات كثيرة ويجلس
الفرق الخام عند سقفه \qty{95}{\%}، بلا أي
معلومة عن الزمن.
\textbf{16.} $10^{-11}\times 10^{9}\times 100\times 2 = 2$ استبدال
بين سلالتين في 100 موقع على مدى مليار سنة. ولانفصال عمره
\qty{10}{Myr} يكون التوقع $0.02$: أي صفر فروق في معظم الأحيان، بلا تمييز.
\textbf{17.} $1200\times 2\times 10^{-9}\times 0.04 = 9.6\times 10^{-8}$
في السنة: فتُتوقَّع أول رامزة وقف بعد نحو \qty{10}{Myr}.
\textbf{18.} مضاعفة الجينوم كله تضاعف كل مورّثة دفعةً واحدة،
فتُحفظ نسب البروتينات المتفاعلة ولا يوجد اختلال جرعة ينتخب
ضد النسخ؛ وأما التضاعف الترادفي المفرد فيغيّر جرعة مورّثة
واحدة وكثيرًا ما يكون ضارًا، أو فائضًا فورًا وحرًا في الاضمحلال.
\textbf{19.} $\mathrm{e}^{-\num{80000}/\num{100000}} = \mathrm{e}^{-0.8} = 0.45$.
\textbf{20.} $\tfrac{2}{3}\times 0.45 = 0.30$: وهو متسق مع
\qty{30}{\%} المرصودة.
\textbf{21.} النصف لكل منهما، أي \qty{15}{\%} من كل الأشجار تضم
الإنسان إلى الغوريلا و\qty{15}{\%} تضم الشمبانزي إلى الغوريلا.
وفائض واضح لأحدهما كان سيدل على تدفق مورّثي بين ذينك النوعين
بعد انفصالهما، وهو ما لا يقدر الانجراف على إنتاجه.
\textbf{22.} عند $N = \num{10000}$: $\mathrm{e}^{-4} = 0.018$، فالمخالفة
\qty{1.2}{\%}. والنسبة \qty{30}{\%} المرصودة تقتضي إذن جمهرة
سلفية من نحو خمسين ألفًا --- أي أضعاف الحجم الفعال
للبشر اليوم.
\textbf{23.} الجمع يفترض أن لكل مورّثة شجرة الأنواع؛ ومع
\omterm{prop:b3:molecular-evolution:ils}{الفرز السلالي غير التام} يكون للمورّثات أشجار كثيرة، وفي
الفروع الداخلية القصيرة قد تختلف أشيع \omterm{prop:b3:molecular-evolution:ils}{شجرة مورّثة} عن شجرة
الأنواع، فتزيد البيانات الأكثر ثقةَ الجواب الخطأ. وأما طريقة
التلاقي فتحسب، لكل شجرة أنواع مرشَّحة، التوزيع المتوقع لأشجار
المورّثات، وتنتقي شجرة الأنواع التي تفسر التواترات المرصودة
أحسن تفسير.
\textbf{24.} \omterm{prop:b3:molecular-evolution:ils}{الفرز السلالي غير التام} وقع في الجمهرة السلفية
المشتركة للبشر المعاصرين جميعًا، فكان سيجعل كل جمهرة بشرية
متساوية القرابة من النياندرتاليين. والفائض عند غير الأفارقة
يعني أن الأليلات النياندرتالية دخلت بعد أن غادر أسلاف غير
الأفارقة أفريقيا: أي اختلاط، قبل نحو خمسين ألف سنة.
\textbf{25.} $\mu \approx 2.3\times 10^{-8}$ لكل موقع في الجيل؛
و$\omega \approx 0.14$؛ ونحو \qty{30}{\%} من أشجار المورّثات
مخالِفة لشجرة الأنواع.
\end{solution}
